\section{Freshness-Aware Project Memory}

The evaluated method treats source freshness and provenance as retrieval
constraints. A task workset may include current project facts, code-bound
evidence, durable memories, optional knowledge or workflow starts, and
verification guidance under a fixed context budget. Code-bound evidence is
classified against the newest complete project snapshot. Evidence that cannot be
verified as current is marked suspect or stale rather than silently presented as
truth.

The evaluation distinguishes two intervention surfaces. The canonical retrieval
track gives every system the same formation input, transfer query, and bounded
context ceiling. It asks about retrieval behavior under equal evidence, not
about a product's unrestricted native session experience. The native \memorix{}
track instead exposes one bounded read-only project-context call and records
whether the agent elects to use it. These surfaces are reported separately.

The artifact also implements a separately labelled native-session formation
diagnostic for Claude Code command hooks. A controller converts private hook
payloads into a portable, snapshot-bound capture, then replays those payloads
through the actual \memorix{} hook command in an isolated checkout and requires
a post-formation storage probe. This path is not an equal-ingestion baseline:
only no-memory and \memorix{} conditions may use it, and its current local
client smoke is excluded from all outcome tables and effect estimates.

The no-memory control keeps the same agent client, actual model route, editable
transfer checkout, ordinary source inspection, verification permissions, timeout,
and filesystem boundary. It lacks only predecessor evidence and a memory tool.
Any condition-specific denial of ordinary current-code work invalidates a run.
